\chapter*{ВВЕДЕНИЕ}
\addcontentsline{toc}{chapter}{ВВЕДЕНИЕ}

% TODO отчитать введение

Трехмерная реконструкция представляет собой процесс создания цифровых 3D-моделей объектов и сцен на основе двумерных изображений.
Это одно из активно развивающихся направлений на стыке компьютерного зрения, графики и машинного обучения.
Ключевым компонентом процесса является построение карты глубины — двумерной структуры данных, где каждый элемент соответствует расстоянию от камеры до объекта в реальном мире.
На основе подобных карт становится возможным восстановление объемной структуры сцены.

Технологии трехмерной реконструкции находят широкое применение в различных областях:
\begin{itemize}[wide=12.5mm, leftmargin=12.5mm]
    \item медицина~\cite{kt, med2, med3},
    \item культура и история~\cite{kult1, kult2},
    \item промышленное производство~\cite{prom1, prom2, prom3},
    \item индустрия развлечений~\cite{play1, play2, play3},
    \item системы компьютерного зрения~\cite{vis1, vis2, vis3}.
\end{itemize}

Работа с архивными фотографиями представляет особый интерес в контексте трехмерной реконструкции по следующим причинам:

\begin{enumerate}
    \item \textbf{Сохранение культурного наследия}: Архивные фотографии зачастую являются единственными свидетельствами существования утраченных культурных объектов, архитектурных сооружений и исторических ландшафтов. Их преобразование в трехмерный формат способствует более полному сохранению исторической памяти.
    
    \item \textbf{Научная и образовательная ценность}: Трехмерные реконструкции исторических объектов открывают новые возможности для исследований и образовательных проектов, делая историческое наследие более доступным и понятным для широкой аудитории.
    
    \item \textbf{Специфика исходных данных}: Архивные фотографии характеризуются различными видами деградации (выцветание, пятна, физические повреждения), неоднородным освещением и отсутствием дополнительной метаинформации, что требует разработки специализированных алгоритмов реконструкции.
\end{enumerate}

Разработанный в рамках данной работы гибридный метод восстановления глубины из монокулярных архивных изображений обладает следующими инновационными характеристиками:

\begin{enumerate}
    \item \textbf{Комбинирование подходов}: Предложенный метод объединяет преимущества нейросетевого подхода (MiDaS), обеспечивающего точное восстановление общей геометрии сцены, и классического метода Shape from Shading (SFS), сохраняющего детализацию и устойчивость к деградации.
    
    \item \textbf{Адаптивность к качеству изображений}: Метод динамически адаптирует стратегию восстановления в зависимости от степени деградации исходного материала, что особенно важно при работе с архивными фотографиями.
    
    \item \textbf{Устойчивость к дефектам}: Разработанное решение демонстрирует значительную устойчивость к типичным дефектам архивных фотографий, сохраняя приемлемую точность реконструкции даже при сильной деградации исходных изображений.
    
    \item \textbf{Отсутствие требований к обучающей выборке}: Метод не требует наличия обучающих данных с эталонной глубиной для архивных изображений, что делает его практически применимым в реальных сценариях работы с историческими коллекциями.
\end{enumerate}

\textbf{Цель работы}: Разработка и исследование гибридного метода восстановления трехмерной информации из архивных фотографий, устойчивого к различным типам деградации исходных изображений.

\textbf{Задачи исследования}:
\begin{enumerate}
    \item Провести анализ существующих методов восстановления глубины из монокулярных изображений и их применимости к архивным фотографиям.
    
    \item Разработать формальную модель гибридного метода, объединяющего нейросетевой подход (MiDaS) и классический метод Shape from Shading.
    
    \item Реализовать программный прототип предложенного метода с использованием современных технологий обработки изображений и машинного обучения.
    
    \item Провести экспериментальное исследование эффективности разработанного метода на изображениях с различной степенью деградации.
    
    \item Сформировать рекомендации по применению разработанного метода в зависимости от характеристик исходных архивных фотографий.
\end{enumerate}

Результаты исследования имеют прикладное значение для задач цифровой сохранности культурного наследия, виртуальных реконструкций исторических объектов и повышения доступности архивных материалов через их представление в современных цифровых форматах.
